\section{Lecture 8 (September 24th)}
\begin{rmk}
(I) The vacuum energy $E$ is the sum of all ground state energies for the harmonic oscillator, and therefore is divergent! We might think of possible resolutions. 
\begin{itemize}
\item[(i)] Ignore, as energy is only meanful as a difference 
\item[(ii)] Introducing some regulator called the IR regulator, taking
\[\delta ^{(3)}(\vec{p})=\int \dfrac{d^3\vec{x}}{(2\pi )^3}\,e^{i\vec{p}\cdot \vec{x}}\longrightarrow \delta ^{(3)}({\bf 0})=\int \dfrac{d^3\vec{x}}{(2\pi )^3}=\dfrac{V}{(2\pi )^3}\]
which implies that
\[E_{0}=\dfrac{V}{(2\pi )^3}=\int d^3\vec{p}\,(E_{p}/2)\]
\item[(iii)] Renormalisation, and redefining our theory by employing a infinite counterterm 
\[\mathcal{L}_{ren}=-\dfrac{1}{2}(\partial \phi )^2-\dfrac{1}{2}m\phi ^2+\mathrm{\Lambda }_{0}\longrightarrow \hat{H}_{rem}=\hat{H}-\int d^3\vec{x}\,\mathrm{\Lambda }_{0}\]
where $\mathrm{\Lambda }_{0}$ would be the cosmological constant and we would set it such that the integral in the right would become $E_0$. In this case, the action would become infinite, but the observables would become finite!
\end{itemize}
(II) Another important remark is that the order of the operators matter. We define the normal ordering to be putting creation operators to the left. Well, how could this be helpful? 

\bigskip 

(III) The Klein-Gordan field operators satisfy the time evolution property in the Heisenberg picture. Well, see that
\[\begin{align*}
e^{i\hat{H}x^{0}}\hat{\phi }(0,\vec{x})e^{-i\hat{H}x^{0}}=&\int \dfrac{d^3\vec{p}}{\sqrt{(2\pi )^32E_{p}}}e^{i\hat{H}x^{0}}\Big(\hat{a}(\vec{p})e^{i\vec{p}\cdot \vec{x}}+\hat{a}(\vec{p})^{\dagger}e^{-i\vec{p}\cdot \vec{x}}\Big)e^{-i\hat{H}x^{0}} \\
=&\int \dfrac{d^3\vec{p}}{\sqrt{(2\pi )^32E_{p}}}\Big(\hat{a}(\vec{p})e^{-iE_{p}x^{0}+i\vec{p}\cdot \vec{x}}+\hat{a}(\vec{p})^{\dagger}e^{iE_{p}x^{0}-i\vec{p}\cdot \vec{x}}\Big)\\
=&\hat{\phi }(x^{0},\vec{x})=\hat{\phi }(x)
\end{align*}\]
with
\[\begin{align*}
[\hat{H},\hat{a}(\vec{p})]=&\int d^3\vec{q}\,E_{q}[\hat{a}_{q}^{\dagger}\hat{a}_{q},\hat{a}_{p}]\\
=&\int d^3\vec{q}\,E_{q}\Big(\hat{a}^{\dagger}_{q}[\hat{a}_{q},\hat{a}_{p}]+[\hat{a}_{p}^{\dagger},\hat{a}_{p}]\hat{a}_{q}\Big)\\
=&-E_{p}\hat{a}(\vec{p})
\end{align*}\]
and
\[[\hat{H},\hat{a}(\vec{p})^{\dagger}]=E_{p}\hat{a}(\vec{p})^{\dagger}\]
Notably, the commutator relations with the ladder operators and the Hamiltonian allows the fields to transform unitarily under the time-evoluation operator. For general spacetime translations, we'll see also that
\[e^{-i\hat{P}_{\mu }a^{\mu }}\hat{\phi }(x)e^{i\hat{P}_{\mu }a^{\mu }}=\hat{\phi }(x+a)\]
this characterises the translation invariance of a scalar field operator.
\end{rmk}
\vspace{2ex}
\begin{rmk}
(Testing the Fock space) Acting the time-ordered Hamiltonian on the momentum ket, we find
\[\hat{H}|\vec{p}\rangle =\int d^3\vec{q}\,E_{q}\hat{a}_{q}^{\dagger}[\hat{a}_{q},\hat{a}_{p}^{\dagger}]|0\rangle=E_{p}\hat{a}_{p}^{\dagger}|0\rangle =E_{p}|\vec{p}\rangle  \]
and
\[\hat{H}|\vec{p}_{1},\ldots ,\vec{p}_{N}\rangle =(E_{p_1}+\ldots E_{p_{n}})|\vec{p}_{1},\ldots ,\vec{p}_{N}\rangle \]
which is the intuition we have. Also,
\[\hat{P}|\vec{p}_{1},\ldots ,\vec{p}_{n}\rangle \]
\end{rmk}
We also make the note that if we deal with complex Klein-Gordan theory, we get another interesting charge which is associated with the $U(1)$ symmetry!
\vspace{2ex}
\begin{defi}
(Feynman propagator and the time ordered product) The Feynman propagator is a crucial ingredient in interacting quantum field theory.
\[\mathrm{\Delta }_{F}(x,y)=\langle 0|T\hat{\phi }(x)\hat{\phi }(y)|0\rangle =\mathrm{\Delta }_{F}(x-y)\]
where the equality followings from translation symmetry. The $T$ that we wrote is the time ordered product defined as
\[T\hat{A}(x)\hat{B}(y)=\mathrm{\Theta}(x^{0}-y^{0})\hat{A}(x)\hat{B}(y)\pm \mathrm{\Theta}(y^{0}-x^{0})\hat{B}(y)\hat{A}(x)\]
where the plus is for bosonic cases and the minus is for fermionic cases. Let's try evaluating the Feynman propagator. Having
\[\hat{\phi }(x)=\int \dfrac{d^3\vec{p}}{\sqrt{(2\pi )^32E_{p}}}\Big(\hat{a}_{p}e^{ip\cdot x}+\hat{a}_{p}^{\dagger}e^{-ip\cdot x}\Big)\qquad\&\qquad \langle 0|\hat{a}_{p}^{\dagger}=\hat{a}_{p}|0\rangle =0\]
This then gives
\[\begin{align*}
\mathrm{\Delta }_{F}(x-y)=&\int \dfrac{d^3\vec{p}d^3\vec{q}}{(2\pi )^3\sqrt{2E_{p}2E_{q}}}\bigg(\mathrm{\Theta}(x^{0}-y^{0})\langle 0|[\hat{a}_{q},\hat{a}_{q}^{\dagger}]|0\rangle e^{i(p\cdot x-q\cdot y)}\\
&+\mathrm{\Theta}(y^{0}-x^{0})\langle 0|[\hat{a}_{q},\hat{a}_{p}^{\dagger}]|0\rangle e^{-i(p\cdot x-q\cdot y)}
\\
=&\int \dfrac{d^3\vec{p}}{(2\pi )^32E_{p}}\Big(\mathrm{\Theta}(x^{0}-y^{0})e^{ip\cdot (x-y)}+\mathrm{\Theta}(y^{0}-x^{0})e^{-ip\cdot (x-y)}\Big)\\
=&\int \dfrac{d^3\vec{p}}{(2\pi )^3}e^{i\vec{p}\cdot (\vec{x}-\vec{y})}\times \dfrac{1}{2E_{p}}\Big(\mathrm{\Theta}(x^{0}-y^{0})e^{-iE_{p}(x^{0}-y^{0})}+\mathrm{\Theta}(y^{0}-x^{0})e^{-iE_{p}(x^{0}-y^{0})}\Big)\\
=&\int \dfrac{d^{4}p}{(2\pi )^{4}}\dfrac{-ie^{ip\cdot (x-y)}}{p^2+m^2-i\varepsilon }
\end{align*}\]
where $\varepsilon \varepsilon_0>0$ is arbitrarily small. Surprisingly,
\[\int \dfrac{dp^{0}}{2\pi i}\dfrac{e^{-ip^{0}(x^{0}-y^{0})}}{-(p^{0})^2+\vec{p}^2+m^2-i\varepsilon }=\int \dfrac{dp^{0}}{2\pi i}\dfrac{-e^{-ip^{0}(x^{0}-y^{0})}}{(p^{0}+E_p-i\varepsilon )(p^{0}-E_{p}+i\varepsilon )}\]
with two residues at $-E_{p}$ and $E_{p}$ each. We find
\[\mathrm{\Theta}(x^{0}-y^{0})\Big(\dfrac{e^{-iE_{p}(x^{0}-y^{0})}}{2E_{p}}\Big)+\mathrm{\Theta}(y^{0}-x^{0})\Big(\dfrac{-e^{iE_{p}(x^{0}-y^{0})}}{-2E_{p}}\Big)\]
\end{defi}
\vspace{2ex}
\begin{prop}
The Feynman propagator acts as a Green's function of the Klein-Gordan equation, with
\[\begin{align*}
(-\partial _{x}^2+m^2-i\varepsilon )\mathrm{\Delta }_{F}(x-y)\\
=&(-\partial _{x}^2+m^2-i\varepsilon )\int \dfrac{d^{4}p}{(2\pi )^{4}}\cdot \dfrac{-ie^{ip\cdot (x-y)}}{p^2+m^2-i\varepsilon }\\
=&\int \dfrac{d^{4}p}{(2\pi )^{4}}\,(p^2+m^2-i\varepsilon )\dfrac{-ie^{ip\cdot (x-y)}}{p^2+m^2-i\varepsilon }=-i\delta ^{(4)}(x-y)
\end{align*}\]
where $p^{\mu }p_{\mu }=p^2$.
\end{prop}
\vspace{2ex}

